%!TEX program = xelatex
%!TEX root = ../all_beamer.tex

%% ============================================================
%%  Foundations of Machine Learning
%%  Chapter 3 – Loss Functions & Training
%%  Beamer Presentation
%%  Giorgio Gambosi – University of Rome "Tor Vergata"
%%  A.Y. 2025–2026
%% ============================================================
%%% !TEX root = all_beamer.tex
% ==============================================================================
% HEADER LATEX PER PRESENTAZIONI BEAMER
% Corso di Machine Learning - Giorgio Gambosi
% Università di Roma "Tor Vergata"
% ==============================================================================

% ------------------------------------------------------------------------------
% CLASSE DOCUMENTO E METADATI
% ------------------------------------------------------------------------------
\documentclass[aspectratio=169, 10pt]{beamer}

\subtitle{Course of Machine Learning}
\author{Giorgio Gambosi}
\institute[Tor Vergata]{Master's Degree in Computer Science\\University of Rome ``Tor Vergata''}
\date{A.Y. 2025--2026}

% ------------------------------------------------------------------------------
% PACCHETTI MATEMATICI
% ------------------------------------------------------------------------------
\usepackage{amsmath, amssymb, amsfonts, mathtools, amsthm}
\usepackage{bm}              % Simboli matematici in grassetto
\usepackage{upgreek}         % Lettere greche dritte
\usepackage{derivative}      % Gestione semplificata delle derivate

% ------------------------------------------------------------------------------
% PACCHETTI GRAFICI E VISUALIZZAZIONE
% ------------------------------------------------------------------------------
\usepackage{graphicx}        % Inclusione immagini
\usepackage{xcolor}          % Gestione colori
\usepackage{xkcdcolors}      % Colori aggiuntivi stile xkcd
\usepackage{microtype}       % Miglioramenti tipografici

% ------------------------------------------------------------------------------
% PACCHETTI PER TABELLE
% ------------------------------------------------------------------------------
\usepackage{booktabs}        % Tabelle professionali
\usepackage{array}           % Estensioni per tabelle
\usepackage{multirow}        % Celle multi-riga
\usepackage{multicol}        % Layout multi-colonna

% ------------------------------------------------------------------------------
% TIKZ E PGFPLOTS
% ------------------------------------------------------------------------------
\usepackage{tikz}
\usepackage{tikz-network}
\usepackage[customcolors]{hf-tikz}
\usepackage{pgfplots}
\pgfplotsset{compat=1.18}
\usepgfplotslibrary{fillbetween}

% Librerie TikZ
\usetikzlibrary{arrows, arrows.meta, decorations.pathmorphing, decorations.pathreplacing}
\usetikzlibrary{decorations.markings, snakes, positioning, backgrounds, fit, shadows}
\usetikzlibrary{automata, patterns, calc, intersections, through, shapes}
\usetikzlibrary{shapes.geometric, shapes.arrows, scopes, chains, matrix}

% ------------------------------------------------------------------------------
% ALTRI PACCHETTI
% ------------------------------------------------------------------------------
\usepackage[skins]{tcolorbox}  % Box colorati
\usepackage{subcaption}         % Sottocaption per figure
\usepackage{listings}           % Inclusione codice sorgente

% ------------------------------------------------------------------------------
% TEMA E STILE BEAMER
% ------------------------------------------------------------------------------
\usetheme{Madrid}
\usecolortheme{whale}
\useinnertheme{rounded}
\usefonttheme{professionalfonts}
\setbeamertemplate{navigation symbols}{}  % Rimuove simboli di navigazione

% ------------------------------------------------------------------------------
% DEFINIZIONE COLORI PERSONALIZZATI
% ------------------------------------------------------------------------------
% Colori principali del tema
\definecolor{MainBlue}{RGB}{31,56,100}      % Blu principale (header, titoli)
\definecolor{AccentBlue}{RGB}{70,130,180}   % Blu accento (blocchi)
\definecolor{AlertRed}{RGB}{160,30,30}      % Rosso per alert
\definecolor{ExGreen}{RGB}{30,120,60}       % Verde per esempi
\definecolor{BlockBg}{RGB}{220,232,246}     % Sfondo blocchi
\definecolor{torgray}{RGB}{220,225,235}     % Grigio Tor Vergata

% Colori alternativi (Tor Vergata style)
\definecolor{TvDark}{RGB}{30,50,80}         % Blu scuro header
\definecolor{TvMid}{RGB}{70,100,140}        % Blu medio blocchi
\definecolor{TvLight}{RGB}{210,220,235}     % Blu chiaro sfondi
\definecolor{TvAlert}{RGB}{160,30,30}       % Rosso alert
\definecolor{TvGreen}{RGB}{20,100,50}       % Verde esempi

% Colori supplementari
\definecolor{torred}{RGB}{153,0,0}
\definecolor{torcyan}{RGB}{0,112,192}
\definecolor{cerise}{HTML}{DA3A62}
\definecolor{teal}{HTML}{029386}
\definecolor{slateblue}{HTML}{5B5EA6}
\definecolor{orange}{HTML}{F97306}

% Colori per nodi TikZ
\definecolor{nodeblue}{RGB}{101,140,196}    % xkcd Faded Blue

% ------------------------------------------------------------------------------
% CONFIGURAZIONE COLORI BEAMER
% ------------------------------------------------------------------------------
\setbeamercolor{frametitle}{bg=MainBlue, fg=white}
\setbeamercolor{title}{bg=MainBlue, fg=white}
\setbeamercolor{structure}{fg=MainBlue}

% Blocchi standard
\setbeamercolor{block title}{bg=AccentBlue, fg=white}
\setbeamercolor{block body}{bg=BlockBg}

% Blocchi alert
\setbeamercolor{block title alerted}{bg=AlertRed, fg=white}
\setbeamercolor{block body alerted}{bg=red!8}

% Blocchi esempio
\setbeamercolor{block title example}{bg=ExGreen, fg=white}
\setbeamercolor{block body example}{bg=green!7}

% ------------------------------------------------------------------------------
% CONFIGURAZIONE FONT BEAMER
% ------------------------------------------------------------------------------
\setbeamerfont{frametitle}{size=\normalsize,series=\bfseries}
\setbeamerfont{block title}{size=\small,series=\bfseries}

% ------------------------------------------------------------------------------
% STILI TIKZ PERSONALIZZATI
% ------------------------------------------------------------------------------
\tikzset{
  >=latex,  % Frecce stile LaTeX di default
  % Nodi per reti neurali
  node/.style={thick,circle,draw=myblue,minimum size=22,inner sep=0.5,outer sep=0.6},
  node in/.style={node,green!20!black,draw=mygreen!30!black,fill=mygreen!25},
  node hidden/.style={node,blue!20!black,draw=myblue!30!black,fill=myblue!20},
  node convol/.style={node,orange!20!black,draw=myorange!30!black,fill=myorange!20},
  node out/.style={node,red!20!black,draw=myred!30!black,fill=myred!20},
  % Connessioni
  connect/.style={thick,mydarkblue},
  connect arrow/.style={-{Latex[length=4,width=3.5]},thick,mydarkblue,shorten <=0.5,shorten >=1},
  % Stili numerati per facile mappatura
  node 1/.style={node in},
  node 2/.style={node hidden},
  node 3/.style={node out}
}

% ------------------------------------------------------------------------------
% AMBIENTI TEOREMI
% ------------------------------------------------------------------------------
\setbeamertemplate{theorems}[numbered]
\newtheorem{mytheorem}{Theorem}
\newtheorem{mydefinition}{Definition}

% ------------------------------------------------------------------------------
% CONFIGURAZIONE LISTINGS (CODICE)
% ------------------------------------------------------------------------------
\lstset{
  language=Python,
  basicstyle=\ttfamily\tiny,
  keywordstyle=\color{MainBlue}\bfseries,
  commentstyle=\color{gray},
  frame=single,
  breaklines=true,
}

% ==============================================================================
% MACRO MATEMATICHE PERSONALIZZATE
% ==============================================================================

% ------------------------------------------------------------------------------
% SPAZI CALLIGRAFICI (Calligraphic spaces)
% ------------------------------------------------------------------------------
\providecommand{\calX}{\mathcal{X}}         % Spazio input
\providecommand{\calY}{\mathcal{Y}}         % Spazio output
\providecommand{\calH}{\mathcal{H}}         % Spazio ipotesi
\providecommand{\calT}{\mathcal{T}}         % Training set
\providecommand{\calA}{\mathcal{A}}         % Algoritmo
\providecommand{\calR}{\mathcal{R}}         % Rischio
\providecommand{\calF}{\mathcal{F}}         % Famiglia di funzioni
\providecommand{\calL}{\mathcal{L}}         % Loss function

% Alias per compatibilità
\providecommand{\Rr}{\mathcal{R}}
\providecommand{\Tr}{\mathcal{T}}
\providecommand{\Hr}{\mathcal{H}}

% ------------------------------------------------------------------------------
% RISCHIO E PROBABILITÀ
% ------------------------------------------------------------------------------
\providecommand{\barR}{\overline{\mathcal{R}}}  % Rischio empirico
\providecommand{\EmpR}{\overline{\mathcal{R}}}  % Rischio empirico (alias)
\providecommand{\Rbar}{\overline{\mathcal{R}}}  % Rischio empirico (alias)
\providecommand{\pM}{p_M}                       % Probabilità modello
\providecommand{\pC}{p_C}                       % Probabilità corretta
\providecommand{\Prob}[2]{\mathbb{P}_{#1}\!\left[#2\right]}  % Probabilità

% ------------------------------------------------------------------------------
% FUNZIONI INDICATRICI
% ------------------------------------------------------------------------------
\providecommand{\indicator}[1]{\mathbf{1}\!\left[#1\right]}  % Funzione indicatrice
\providecommand{\ind}[1]{\mathbf{1}[#1]}                      % Versione breve

% ------------------------------------------------------------------------------
% OPERATORI MATEMATICI
% ------------------------------------------------------------------------------
\providecommand{\argmax}{\operatorname*{arg\,max}}  % Argomento del massimo
\providecommand{\argmin}{\operatorname*{arg\,min}}  % Argomento del minimo
\providecommand{\vcd}[1]{\operatorname{VCdim}(#1)} % VC dimension
\providecommand{\defeq}{\stackrel{\text{def}}{=}}  % Definito come

% ------------------------------------------------------------------------------
% VETTORI E MATRICI (Bold Math)
% ------------------------------------------------------------------------------
% Vettori minuscoli
\providecommand{\bx}{\mathbf{x}}            % Vettore x
\providecommand{\bz}{\mathbf{z}}            % Vettore z
\providecommand{\bw}{\mathbf{w}}            % Vettore peso
\providecommand{\btt}{\mathbf{t}}           % Vettore target
\providecommand{\wvec}{\mathbf{w}}          % Vettore peso (alias)
\providecommand{\bth}{\boldsymbol{\theta}}  % Parametro theta

% Matrici maiuscole
\providecommand{\bX}{\mathbf{X}}            % Matrice dati
\providecommand{\bZ}{\mathbf{Z}}            % Matrice Z
\providecommand{\Xmat}{\mathbf{X}}          % Matrice dati (alias)
\providecommand{\Hmat}{\mathbf{H}}          % Matrice H
\providecommand{\Smat}{\mathbf{S}}          % Matrice scatter generica
\providecommand{\bS}{\mathbf{S}}            % Matrice scatter
\providecommand{\Ii}{\mathbf{I}}            % Matrice identità
\providecommand{\bZero}{\mathbf{0}}         % Vettore zero
\providecommand{\bPsi}{\boldsymbol{\Psi}}   % Matrice Psi

% Medie (overline)
\providecommand{\ox}{\overline{\mathbf{x}}} % Media vettore x
\providecommand{\ow}{\overline{\mathbf{w}}} % Media vettore w
\providecommand{\oX}{\overline{\mathbf{X}}} % Media matrice X

% Matrici scatter specifiche (analisi discriminante)
\providecommand{\SW}{\mathbf{S}_W}          % Within-class scatter
\providecommand{\SB}{\mathbf{S}_B}          % Between-class scatter
\providecommand{\ST}{\mathbf{S}_T}          % Total scatter
\providecommand{\GX}{G_{\mathbf{X}}}        % Matrice Gram

% ------------------------------------------------------------------------------
% COMANDI DI FORMATTAZIONE TESTO
% ------------------------------------------------------------------------------
\providecommand{\term}[1]{\textbf{\color{accent}#1}}        % Termine importante
\providecommand{\halert}[1]{{\color{AlertRed}\textbf{#1}}}  % Alert rosso
\providecommand{\mlalert}[1]{{\color{AccentBlue}\textbf{#1}}}  % Alert blu
\providecommand{\mlemph}[1]{{\color{MainBlue}\textit{#1}}}  % Enfasi blu

% ==============================================================================
% FINE HEADER
% ==============================================================================

\newcommand{\blankpage}{
\clearpage
\thispagestyle{empty}
\null
\clearpage
}

%\documentclass[aspectratio=169, 10pt]{beamer}
%
%% ---------- Theme & colours ----------
%\usetheme{Madrid}
%\usecolortheme{whale}
%\useinnertheme{rounded}
%\setbeamertemplate{navigation symbols}{}
%
%\definecolor{MainBlue}{RGB}{31,56,100}
%\definecolor{AccentBlue}{RGB}{70,130,180}
%\definecolor{AlertRed}{RGB}{160,30,30}
%\definecolor{ExGreen}{RGB}{30,120,60}
%\definecolor{BlockBg}{RGB}{220,232,246}
%\definecolor{torgray}{RGB}{220,225,235}
%
%\setbeamercolor{frametitle}{bg=MainBlue, fg=white}
%\setbeamercolor{title}{bg=MainBlue, fg=white}
%\setbeamercolor{block title}{bg=AccentBlue, fg=white}
%\setbeamercolor{block body}{bg=BlockBg}
%\setbeamercolor{block title alerted}{bg=AlertRed, fg=white}
%\setbeamercolor{block body alerted}{bg=red!8}
%\setbeamercolor{block title example}{bg=ExGreen, fg=white}
%\setbeamercolor{block body example}{bg=green!7}
%\setbeamercolor{structure}{fg=MainBlue}
%
%% ---------- Fonts ----------
%\usefonttheme{professionalfonts}
%\setbeamerfont{frametitle}{size=\normalsize,series=\bfseries}
%\setbeamerfont{block title}{size=\small,series=\bfseries}

%% !TEX root = all_beamer.tex
% ==============================================================================
% HEADER LATEX PER PRESENTAZIONI BEAMER
% Corso di Machine Learning - Giorgio Gambosi
% Università di Roma "Tor Vergata"
% ==============================================================================

% ------------------------------------------------------------------------------
% CLASSE DOCUMENTO E METADATI
% ------------------------------------------------------------------------------
\documentclass[aspectratio=169, 10pt]{beamer}

\subtitle{Course of Machine Learning}
\author{Giorgio Gambosi}
\institute[Tor Vergata]{Master's Degree in Computer Science\\University of Rome ``Tor Vergata''}
\date{A.Y. 2025--2026}

% ------------------------------------------------------------------------------
% PACCHETTI MATEMATICI
% ------------------------------------------------------------------------------
\usepackage{amsmath, amssymb, amsfonts, mathtools, amsthm}
\usepackage{bm}              % Simboli matematici in grassetto
\usepackage{upgreek}         % Lettere greche dritte
\usepackage{derivative}      % Gestione semplificata delle derivate

% ------------------------------------------------------------------------------
% PACCHETTI GRAFICI E VISUALIZZAZIONE
% ------------------------------------------------------------------------------
\usepackage{graphicx}        % Inclusione immagini
\usepackage{xcolor}          % Gestione colori
\usepackage{xkcdcolors}      % Colori aggiuntivi stile xkcd
\usepackage{microtype}       % Miglioramenti tipografici

% ------------------------------------------------------------------------------
% PACCHETTI PER TABELLE
% ------------------------------------------------------------------------------
\usepackage{booktabs}        % Tabelle professionali
\usepackage{array}           % Estensioni per tabelle
\usepackage{multirow}        % Celle multi-riga
\usepackage{multicol}        % Layout multi-colonna

% ------------------------------------------------------------------------------
% TIKZ E PGFPLOTS
% ------------------------------------------------------------------------------
\usepackage{tikz}
\usepackage{tikz-network}
\usepackage[customcolors]{hf-tikz}
\usepackage{pgfplots}
\pgfplotsset{compat=1.18}
\usepgfplotslibrary{fillbetween}

% Librerie TikZ
\usetikzlibrary{arrows, arrows.meta, decorations.pathmorphing, decorations.pathreplacing}
\usetikzlibrary{decorations.markings, snakes, positioning, backgrounds, fit, shadows}
\usetikzlibrary{automata, patterns, calc, intersections, through, shapes}
\usetikzlibrary{shapes.geometric, shapes.arrows, scopes, chains, matrix}

% ------------------------------------------------------------------------------
% ALTRI PACCHETTI
% ------------------------------------------------------------------------------
\usepackage[skins]{tcolorbox}  % Box colorati
\usepackage{subcaption}         % Sottocaption per figure
\usepackage{listings}           % Inclusione codice sorgente

% ------------------------------------------------------------------------------
% TEMA E STILE BEAMER
% ------------------------------------------------------------------------------
\usetheme{Madrid}
\usecolortheme{whale}
\useinnertheme{rounded}
\usefonttheme{professionalfonts}
\setbeamertemplate{navigation symbols}{}  % Rimuove simboli di navigazione

% ------------------------------------------------------------------------------
% DEFINIZIONE COLORI PERSONALIZZATI
% ------------------------------------------------------------------------------
% Colori principali del tema
\definecolor{MainBlue}{RGB}{31,56,100}      % Blu principale (header, titoli)
\definecolor{AccentBlue}{RGB}{70,130,180}   % Blu accento (blocchi)
\definecolor{AlertRed}{RGB}{160,30,30}      % Rosso per alert
\definecolor{ExGreen}{RGB}{30,120,60}       % Verde per esempi
\definecolor{BlockBg}{RGB}{220,232,246}     % Sfondo blocchi
\definecolor{torgray}{RGB}{220,225,235}     % Grigio Tor Vergata

% Colori alternativi (Tor Vergata style)
\definecolor{TvDark}{RGB}{30,50,80}         % Blu scuro header
\definecolor{TvMid}{RGB}{70,100,140}        % Blu medio blocchi
\definecolor{TvLight}{RGB}{210,220,235}     % Blu chiaro sfondi
\definecolor{TvAlert}{RGB}{160,30,30}       % Rosso alert
\definecolor{TvGreen}{RGB}{20,100,50}       % Verde esempi

% Colori supplementari
\definecolor{torred}{RGB}{153,0,0}
\definecolor{torcyan}{RGB}{0,112,192}
\definecolor{cerise}{HTML}{DA3A62}
\definecolor{teal}{HTML}{029386}
\definecolor{slateblue}{HTML}{5B5EA6}
\definecolor{orange}{HTML}{F97306}

% Colori per nodi TikZ
\definecolor{nodeblue}{RGB}{101,140,196}    % xkcd Faded Blue

% ------------------------------------------------------------------------------
% CONFIGURAZIONE COLORI BEAMER
% ------------------------------------------------------------------------------
\setbeamercolor{frametitle}{bg=MainBlue, fg=white}
\setbeamercolor{title}{bg=MainBlue, fg=white}
\setbeamercolor{structure}{fg=MainBlue}

% Blocchi standard
\setbeamercolor{block title}{bg=AccentBlue, fg=white}
\setbeamercolor{block body}{bg=BlockBg}

% Blocchi alert
\setbeamercolor{block title alerted}{bg=AlertRed, fg=white}
\setbeamercolor{block body alerted}{bg=red!8}

% Blocchi esempio
\setbeamercolor{block title example}{bg=ExGreen, fg=white}
\setbeamercolor{block body example}{bg=green!7}

% ------------------------------------------------------------------------------
% CONFIGURAZIONE FONT BEAMER
% ------------------------------------------------------------------------------
\setbeamerfont{frametitle}{size=\normalsize,series=\bfseries}
\setbeamerfont{block title}{size=\small,series=\bfseries}

% ------------------------------------------------------------------------------
% STILI TIKZ PERSONALIZZATI
% ------------------------------------------------------------------------------
\tikzset{
  >=latex,  % Frecce stile LaTeX di default
  % Nodi per reti neurali
  node/.style={thick,circle,draw=myblue,minimum size=22,inner sep=0.5,outer sep=0.6},
  node in/.style={node,green!20!black,draw=mygreen!30!black,fill=mygreen!25},
  node hidden/.style={node,blue!20!black,draw=myblue!30!black,fill=myblue!20},
  node convol/.style={node,orange!20!black,draw=myorange!30!black,fill=myorange!20},
  node out/.style={node,red!20!black,draw=myred!30!black,fill=myred!20},
  % Connessioni
  connect/.style={thick,mydarkblue},
  connect arrow/.style={-{Latex[length=4,width=3.5]},thick,mydarkblue,shorten <=0.5,shorten >=1},
  % Stili numerati per facile mappatura
  node 1/.style={node in},
  node 2/.style={node hidden},
  node 3/.style={node out}
}

% ------------------------------------------------------------------------------
% AMBIENTI TEOREMI
% ------------------------------------------------------------------------------
\setbeamertemplate{theorems}[numbered]
\newtheorem{mytheorem}{Theorem}
\newtheorem{mydefinition}{Definition}

% ------------------------------------------------------------------------------
% CONFIGURAZIONE LISTINGS (CODICE)
% ------------------------------------------------------------------------------
\lstset{
  language=Python,
  basicstyle=\ttfamily\tiny,
  keywordstyle=\color{MainBlue}\bfseries,
  commentstyle=\color{gray},
  frame=single,
  breaklines=true,
}

% ==============================================================================
% MACRO MATEMATICHE PERSONALIZZATE
% ==============================================================================

% ------------------------------------------------------------------------------
% SPAZI CALLIGRAFICI (Calligraphic spaces)
% ------------------------------------------------------------------------------
\providecommand{\calX}{\mathcal{X}}         % Spazio input
\providecommand{\calY}{\mathcal{Y}}         % Spazio output
\providecommand{\calH}{\mathcal{H}}         % Spazio ipotesi
\providecommand{\calT}{\mathcal{T}}         % Training set
\providecommand{\calA}{\mathcal{A}}         % Algoritmo
\providecommand{\calR}{\mathcal{R}}         % Rischio
\providecommand{\calF}{\mathcal{F}}         % Famiglia di funzioni
\providecommand{\calL}{\mathcal{L}}         % Loss function

% Alias per compatibilità
\providecommand{\Rr}{\mathcal{R}}
\providecommand{\Tr}{\mathcal{T}}
\providecommand{\Hr}{\mathcal{H}}

% ------------------------------------------------------------------------------
% RISCHIO E PROBABILITÀ
% ------------------------------------------------------------------------------
\providecommand{\barR}{\overline{\mathcal{R}}}  % Rischio empirico
\providecommand{\EmpR}{\overline{\mathcal{R}}}  % Rischio empirico (alias)
\providecommand{\Rbar}{\overline{\mathcal{R}}}  % Rischio empirico (alias)
\providecommand{\pM}{p_M}                       % Probabilità modello
\providecommand{\pC}{p_C}                       % Probabilità corretta
\providecommand{\Prob}[2]{\mathbb{P}_{#1}\!\left[#2\right]}  % Probabilità

% ------------------------------------------------------------------------------
% FUNZIONI INDICATRICI
% ------------------------------------------------------------------------------
\providecommand{\indicator}[1]{\mathbf{1}\!\left[#1\right]}  % Funzione indicatrice
\providecommand{\ind}[1]{\mathbf{1}[#1]}                      % Versione breve

% ------------------------------------------------------------------------------
% OPERATORI MATEMATICI
% ------------------------------------------------------------------------------
\providecommand{\argmax}{\operatorname*{arg\,max}}  % Argomento del massimo
\providecommand{\argmin}{\operatorname*{arg\,min}}  % Argomento del minimo
\providecommand{\vcd}[1]{\operatorname{VCdim}(#1)} % VC dimension
\providecommand{\defeq}{\stackrel{\text{def}}{=}}  % Definito come

% ------------------------------------------------------------------------------
% VETTORI E MATRICI (Bold Math)
% ------------------------------------------------------------------------------
% Vettori minuscoli
\providecommand{\bx}{\mathbf{x}}            % Vettore x
\providecommand{\bz}{\mathbf{z}}            % Vettore z
\providecommand{\bw}{\mathbf{w}}            % Vettore peso
\providecommand{\btt}{\mathbf{t}}           % Vettore target
\providecommand{\wvec}{\mathbf{w}}          % Vettore peso (alias)
\providecommand{\bth}{\boldsymbol{\theta}}  % Parametro theta

% Matrici maiuscole
\providecommand{\bX}{\mathbf{X}}            % Matrice dati
\providecommand{\bZ}{\mathbf{Z}}            % Matrice Z
\providecommand{\Xmat}{\mathbf{X}}          % Matrice dati (alias)
\providecommand{\Hmat}{\mathbf{H}}          % Matrice H
\providecommand{\Smat}{\mathbf{S}}          % Matrice scatter generica
\providecommand{\bS}{\mathbf{S}}            % Matrice scatter
\providecommand{\Ii}{\mathbf{I}}            % Matrice identità
\providecommand{\bZero}{\mathbf{0}}         % Vettore zero
\providecommand{\bPsi}{\boldsymbol{\Psi}}   % Matrice Psi

% Medie (overline)
\providecommand{\ox}{\overline{\mathbf{x}}} % Media vettore x
\providecommand{\ow}{\overline{\mathbf{w}}} % Media vettore w
\providecommand{\oX}{\overline{\mathbf{X}}} % Media matrice X

% Matrici scatter specifiche (analisi discriminante)
\providecommand{\SW}{\mathbf{S}_W}          % Within-class scatter
\providecommand{\SB}{\mathbf{S}_B}          % Between-class scatter
\providecommand{\ST}{\mathbf{S}_T}          % Total scatter
\providecommand{\GX}{G_{\mathbf{X}}}        % Matrice Gram

% ------------------------------------------------------------------------------
% COMANDI DI FORMATTAZIONE TESTO
% ------------------------------------------------------------------------------
\providecommand{\term}[1]{\textbf{\color{accent}#1}}        % Termine importante
\providecommand{\halert}[1]{{\color{AlertRed}\textbf{#1}}}  % Alert rosso
\providecommand{\mlalert}[1]{{\color{AccentBlue}\textbf{#1}}}  % Alert blu
\providecommand{\mlemph}[1]{{\color{MainBlue}\textit{#1}}}  % Enfasi blu

% ==============================================================================
% FINE HEADER
% ==============================================================================

\newcommand{\blankpage}{
\clearpage
\thispagestyle{empty}
\null
\clearpage
}

% ---------- Packages ----------
%\usepackage{amsmath,amssymb,bm}
%\usepackage{mathtools}
%\usepackage{pgfplots}
%\pgfplotsset{compat=1.18}
%\usepackage{tikz}
%\usetikzlibrary{arrows.meta,shapes,positioning}
%\usepackage{booktabs}
%% indicator function via mathbb
%\usepackage{microtype}
%\usepackage{xcolor}
%\usepackage{graphicx}

% ---------- Custom colours for plots ----------
%\definecolor{cerise}{HTML}{DA3A62}
%\definecolor{teal}{HTML}{029386}
%\definecolor{slateblue}{HTML}{5B5EA6}
%\definecolor{orange}{HTML}{F97306}
%
%% ---------- Macros ----------
%
%\newcommand{\calL}{\mathcal{L}}
%\newcommand{\calR}{\mathcal{R}}
%\newcommand{\calT}{\mathcal{T}}
%\newcommand{\calY}{\mathcal{Y}}
%\newcommand{\calX}{\mathcal{X}}
%\newcommand{\calH}{\mathcal{H}}
%\newcommand{\EmpR}{\overline{\mathcal{R}}}

% ---------- Beamer settings ----------
%\setbeamertemplate{navigation symbols}{}
%\setbeamertemplate{footline}{%
%  \leavevmode%
%  \hbox{%
%    \begin{beamercolorbox}[wd=.4\paperwidth,ht=2.25ex,dp=1ex,left]{section in head/foot}%
%      \hspace{1em}\insertshortauthor
%    \end{beamercolorbox}%
%    \begin{beamercolorbox}[wd=.4\paperwidth,ht=2.25ex,dp=1ex,center]{section in head/foot}%
%      Foundations of Machine Learning
%    \end{beamercolorbox}%
%    \begin{beamercolorbox}[wd=.2\paperwidth,ht=2.25ex,dp=1ex,right]{section in head/foot}%
%      A.Y.\ 2025–2026\hspace{0.5em}\insertframenumber\,/\,\inserttotalframenumber\hspace{0.5em}
%    \end{beamercolorbox}}%
%}
%
%\setbeamertemplate{itemize item}{\textbullet}
%\setbeamertemplate{itemize subitem}{--}

% ---------- Title ----------
\title[Loss functions]{Loss functions}
%\subtitle{Course of Machine Learning}
%\author{Giorgio Gambosi}
%\institute{Master's Degree in Computer Science\\University of Rome ``Tor Vergata''}
%\date{A.Y.\ 2025–2026}
%
%% ============================================================
%\begin{document}
%% ============================================================

% --- Title frame ---
\begin{frame}
  \maketitle
\end{frame}

% --- Outline ---
\begin{frame}{Outline}
  \tableofcontents
\end{frame}

% ============================================================
\section{Loss Functions}
% ============================================================

% ---- 1 ----
\begin{frame}{The Loss Function — Concept}
  \begin{block}{Definition}
    The \textbf{loss function} $\mathcal{L}:\mathcal{Y}\times\mathcal{Y}\to\Real$
    measures the \alert{cost} of predicting $y$ when the true target is $t$.
  \end{block}

  \vspace{.5em}
  \begin{itemize}
    \item Both $y$ (prediction) and $t$ (target) belong to the output space $\calY$.
    \item $\calY$ can be: real numbers (regression), class labels (classification), structured outputs.
    \item For a single input $\x$ with true output $t$, the pointwise quality is:
      \[
        \calR(\x, t) = L\!\left(h(\x),\, t\right)
      \]
    \item The loss function drives the \emph{definition} of what "good predictions" means.
  \end{itemize}
\end{frame}

% ---- 2 ----
\begin{frame}{Empirical Risk}
  \begin{block}{Definition}
    The \textbf{empirical risk} aggregates the loss over the training set $\calT$:
    \[
      \EmpR_{\calT}(h)
      = \frac{1}{|\calT|}\sum_{(\x,t)\in\calT} L\!\left(h(\x), t\right)
    \]
  \end{block}

  \vspace{.5em}
  \begin{itemize}
    \item Measures the \emph{average} error of $h$ on the available data.
    \item Lower value $\Rightarrow$ better fit to training data.
    \item \alert{Caveat}: minimising empirical risk $\neq$ minimising generalisation error.
    \item True goal: perform well on \emph{unseen} data (generalisation).
  \end{itemize}
\end{frame}

% ---- 3 ----
\begin{frame}{Parametric Models \& Overall Loss}
  \begin{itemize}
    \item In practice, $h = h_{\Vtheta}$ where $\Vtheta = (\theta_0,\ldots,\theta_d)$ are learned parameters.
    \item Training = searching for the $\Vtheta$ minimising:
      \[
        \calL(\Vtheta;\calT) = \sum_{i=1}^{n} L_i(\Vtheta)
        \quad \text{where}\quad L_i = L(\Vtheta;\x_i,t_i)
      \]
    \item The overall loss sums individual contributions from each example.
  \end{itemize}

  \vspace{.5em}
  \begin{exampleblock}{Interpretation}
    Each data point $(\x_i, t_i)$ contributes a term $L_i(\Vtheta)$ that penalises
    the model for predicting incorrectly on that specific example.
    Training adjusts $\Vtheta$ to reduce this total penalty.
  \end{exampleblock}
\end{frame}

% ============================================================
\section{Optimisation of the Loss}
% ============================================================

% ---- 4 ----
\begin{frame}{Goal: Finding a Global Minimum}
  \begin{itemize}
    \item We seek a \alert{global minimum} of $\calL(\Vtheta;\calT)$ over the parameter space.
    \item A global minimum is a point where the loss is \emph{lower than at every other point}.
  \end{itemize}

  \vspace{.4em}
  \begin{block}{Calculus Approach}
    At any minimum (local or global), the gradient must vanish:
    \[
      \nabla \calL(\Vtheta;\calT) = \Zero
    \]
    i.e., each partial derivative satisfies:
    \[
      \frac{\partial}{\partial \theta_i}\calL(\Vtheta;\calT) = 0 \quad \forall\, i
    \]
  \end{block}

  \begin{itemize}
    \item \textbf{Gradient} $\nabla f$ at a point: vector of steepest ascent.
    \item \textbf{Negative gradient} $-\nabla f$: direction of steepest descent.
  \end{itemize}
\end{frame}

% ---- 5 ----
\begin{frame}{Challenges of Analytical Solutions}
  Setting $\nabla \calL = \Zero$ has several \alert{critical limitations}:

  \vspace{.4em}
  \begin{itemize}
    \item \textbf{Multiple solutions}: critical points include minima, maxima, \emph{and} saddle points.
    \item \textbf{Analytical intractability}: complex models (e.g.\ neural networks) rarely admit closed-form solutions.
    \item \textbf{Non-linearity}: deep networks compose many non-linear functions — no algebraic formula exists.
  \end{itemize}

  \vspace{.5em}
  \begin{alertblock}{Consequence}
    In most practical settings we must resort to \textbf{iterative numerical methods}.
  \end{alertblock}
\end{frame}

% ============================================================
\section{Gradient Descent}
% ============================================================

% ---- 6 ----
\begin{frame}{Gradient Descent — Principle}
  \begin{block}{Core idea}
    Move iteratively in the direction of \alert{steepest descent} of the empirical risk.
  \end{block}

  \vspace{.4em}
  \begin{itemize}
    \item Initialise parameters: $\Vtheta^{(0)} = (\theta_0^{(0)}, \ldots, \theta_d^{(0)})$ (often random).
    \item At each iteration $i$, update \emph{every} parameter:
      \[
        \theta_j^{(i)} \leftarrow \theta_j^{(i-1)}
        - \eta\,\frac{\partial}{\partial \theta_j}
          \EmpR_{\calT}(\Vtheta)\Big|_{\Vtheta^{(i-1)}}
      \]
    \item $\eta > 0$ is the \textbf{learning rate}: controls the step size.
  \end{itemize}
\end{frame}

% ---- 7 ----
\begin{frame}{Gradient Descent — Matrix Form}
  In compact matrix notation, the update becomes:
  \[
    \Vtheta^{(i)} \leftarrow \Vtheta^{(i-1)}
    - \frac{\eta}{|\calT|}\sum_{(\x,t)\in\calT}
      \nabla L\!\left(h_{\Vtheta}(\x), t\right)\Big|_{\Vtheta^{(i-1)}}
  \]

  \vspace{.4em}
  \begin{itemize}
    \item All parameters are updated \emph{simultaneously}.
    \item The gradient is evaluated on the \emph{entire} training set at each step.
    \item This variant is called \textbf{batch gradient descent}.
  \end{itemize}

  \vspace{.4em}
  \begin{exampleblock}{Intuition}
    Think of a ball rolling down a landscape: at each step it moves in the direction
    the slope is steepest, eventually settling in a valley.
  \end{exampleblock}
\end{frame}

% ---- 8 ----
\begin{frame}{The Learning Rate $\eta$}
  \begin{columns}[T]
    \column{.5\textwidth}
    \begin{block}{Too small $\eta$}
      \begin{itemize}
        \item Very slow convergence.
        \item Reliable but impractical for large models.
      \end{itemize}
    \end{block}
    \vspace{.3em}
    \begin{block}{Too large $\eta$}
      \begin{itemize}
        \item May overshoot the minimum.
        \item Can oscillate or diverge.
      \end{itemize}
    \end{block}

    \column{.5\textwidth}
    \begin{block}{Challenges in practice}
      \begin{itemize}
        \item Many local minima in non-convex landscapes.
        \item Saddle points: gradient $= \Zero$, yet not a minimum.
        \item Convergence speed depends on curvature and $\eta$.
        \item Modern solutions: momentum, adaptive rates (Adam, RMSProp).
      \end{itemize}
    \end{block}
  \end{columns}
\end{frame}

% ============================================================
\section{Convexity}
% ============================================================

% ---- 9 ----
\begin{frame}{Convex Sets}
  \begin{block}{Definition}
    A set $S\subset\Real^d$ is \textbf{convex} if for any $\x_1,\x_2\in S$ and $\lambda\in(0,1)$:
    \[
      \lambda\x_1 + (1-\lambda)\x_2 \in S
    \]
  \end{block}

  \vspace{.3em}
  \begin{itemize}
    \item Any point on the \emph{segment} joining $\x_1$ and $\x_2$ must lie in $S$.
    \item Convex sets have no holes or indentations.
  \end{itemize}

  \vspace{.3em}
  \begin{center}
  \begin{tikzpicture}[scale=.45]
    \draw [fill=torgray] plot [mark=none, smooth cycle] coordinates
      {(-6.5,4) (-2.5,2.5) (1,1.7) (3,1.591) (7.312,3)
       (6.7,7.668) (4.264,10.821) (0.948,10.904) (-5.973,5.294)};
    \coordinate (g) at (3,7);
    \fill (g) circle (2.5pt) node [right] {\small$\x_2$};
    \coordinate (p) at (-3,5);
    \fill (p) circle (2.5pt) node [below] {\small$\x_1$};
    \draw [dashed] (p) -- (g);
    \coordinate (e) at (1,6.33);
    \fill (e) circle (2.5pt)
      node [below right] {\tiny$\lambda\x_1{+}(1{-}\lambda)\x_2$};
  \end{tikzpicture}
  \end{center}
\end{frame}

% ---- 10 ----
\begin{frame}{Convex Functions}
  \begin{block}{Definition}
    $f:\Real^d\to\Real$ is \textbf{convex} if for all $\x_1,\x_2$ and $\lambda\in(0,1)$:
    \[
      f\!\left(\lambda\x_1+(1-\lambda)\x_2\right) \leq
      \lambda f(\x_1)+(1-\lambda) f(\x_2)
    \]
    It is \textbf{strictly convex} if the inequality is strict for $\x_1\neq\x_2$.
  \end{block}

  \vspace{.3em}
  \begin{itemize}
    \item Geometric view: the chord between any two graph points lies \emph{above} the graph.
  \end{itemize}

  \vspace{.3em}
  \begin{center}
  \begin{tikzpicture}[scale=0.7]
    \begin{axis}[
      axis x line=bottom, axis y line=left,
      ticks=none, xmin=-2.2, xmax=1.2, ymin=-.5, ymax=3.5,
      width=6cm, height=4cm]
      \addplot [cerise, thick, domain=-2:1, samples=101]{x^2};
      \addplot [slateblue, thick, domain=-2.1:1, samples=101]{-x+.75};
      \addplot [dashed, domain=-1.5:.5, samples=2] coordinates {(-1.5,2.25)(.5,0.25)};
    \end{axis}
  \end{tikzpicture}
  \end{center}
\end{frame}

% ---- 11 ----
\begin{frame}{Why Convexity Matters}
  \begin{alertblock}{Key theorem}
    If $f$ is \textbf{convex}, any local minimum is also a \textbf{global minimum}.\\
    If $f$ is \textbf{strictly convex}, the global minimum is \textbf{unique}.
  \end{alertblock}

  \vspace{.5em}
  \begin{itemize}
    \item For a convex loss, solving $\nabla\calL(\Vtheta;\calT)=\Zero$ \emph{guarantees} a global minimum.
    \item Gradient descent will converge to the optimum regardless of initialisation.
    \item \textbf{Sum rule}: a (positively) weighted sum of convex functions is convex.
    \item As a consequence, if each $L(\Vtheta;\x_i,t_i)$ is convex, so is $\EmpR_{\calT}{(h_{\Vtheta})}$.
  \end{itemize}

  \vspace{.5em}
  \begin{exampleblock}{Practical implication}
    Designing \emph{convex} loss functions makes optimisation reliable and predictable.
    Non-convex losses (e.g.\ in deep networks) require more sophisticated strategies.
  \end{exampleblock}
\end{frame}

% ---- 12 ----
\begin{frame}{Convexity in Practice}
  \begin{columns}[T]
    \column{.5\textwidth}
    \begin{block}{Favourable cases}
      \begin{itemize}
        \item Linear regression + quadratic loss: \emph{strictly convex}, closed-form solution.
        \item Logistic regression + log loss: \emph{strictly convex}.
        \item Linear SVM + hinge loss: \emph{convex}.
      \end{itemize}
    \end{block}

    \column{.5\textwidth}
    \begin{block}{Unfavourable cases}
      \begin{itemize}
        \item Deep neural networks: \emph{highly non-convex}.
        \item Many saddle points; many local minima.
        \item Training requires careful initialisation and adaptive optimisers.
      \end{itemize}
    \end{block}
  \end{columns}

  \vspace{.5em}
  \begin{itemize}
    \item When convexity cannot be guaranteed, we use advanced methods:
      momentum, second-order approximations, learning-rate schedules, etc.
  \end{itemize}
\end{frame}

% ============================================================
\section{Common Loss Functions for Regression}
% ============================================================

% ---- 13 ----
\begin{frame}{Overview: Loss Functions for Regression}
  \begin{itemize}
    \item In regression, both $t$ and $h(\x)$ are \emph{real-valued}.
    \item Error = some distance between prediction and target.
    \item Different distances $\Rightarrow$ different loss functions with different properties.
  \end{itemize}

  \vspace{.5em}
  \begin{center}
  \begin{tabular}{lll}
    \toprule
    \textbf{Loss} & \textbf{Formula} & \textbf{Property} \\
    \midrule
    Quadratic ($L_2$) & $(y-t)^2$ & Smooth, sensitive to outliers \\
    Absolute ($L_1$)  & $|y-t|$   & Robust, non-smooth at 0 \\
    Huber             & piecewise & Best of both worlds \\
    \bottomrule
  \end{tabular}
  \end{center}
\end{frame}

% ---- 14 ----
\begin{frame}{Quadratic Loss (Squared Error)}
  \begin{block}{Definition}
    \[
      L(y,t) = (y-t)^2
    \]
  \end{block}

  \begin{itemize}
    \item Errors penalised \emph{quadratically}: doubling the error quadruples the loss.
    \item Empirical risk $\Rightarrow$ \textbf{Mean Squared Error} (MSE):
      \[
        \EmpR_{\calT}(h)=\frac{1}{|\calT|}\sum_{(\x,t)\in\calT}(h(\x)-t)^2
      \]
    \item For \textbf{linear regression} $h(\x)=\w^T\x+b$:
      \begin{itemize}
        \item Loss is \emph{strictly convex} in $(\w,b)$.
        \item Gradient is \emph{linear} $\Rightarrow$ closed-form solution (normal equations).
      \end{itemize}
    \item \alert{Weakness}: very sensitive to outliers.
  \end{itemize}
\end{frame}

% ---- 15 ----
\begin{frame}{Quadratic Loss — Plot}
  \begin{center}
  \begin{tikzpicture}[scale=0.85]
    \begin{axis}[
      legend style={at={(1.5,1)},anchor=north east,draw=none,fill=none},
      ytick={0}, yticklabels={$0$},
      xticklabels={,,},
      axis x line=bottom, axis y line=left,
      xlabel={$y$},
      x label style={at={(axis description cs:0.99,0)},anchor=north},
      xmin=-2, xmax=2, ymin=-1.1, ymax=3.5,
      width=8cm, height=5cm]
      \addplot [cerise, thick, domain=-2:2, samples=201]{x^2};
      \addplot [slateblue, thick, domain=-2:2, samples=201]{2*x};
      \addplot [dotted, domain=-3:3, samples=2] coordinates {(-3,0)(3,0)};
      \legend{loss, loss gradient}
    \end{axis}
  \end{tikzpicture}
  \end{center}
  \vspace{-0.3em}
  {\small Quadratic loss function and its gradient.
  Note the parabolic shape and the linear gradient.}
\end{frame}

% ---- 16 ----
\begin{frame}{Quadratic Loss — Gradients for Linear Regression}
  For $h(\x) = \w^T\x + b$, the gradients are:
  \[
    \frac{\partial}{\partial w_i}\calL =
    \sum_{(\x,t)\in\calT}(\w^T\x+b-t)\,x_i
    \qquad
    \frac{\partial}{\partial b}\calL =
    \sum_{(\x,t)\in\calT}(\w^T\x+b-t)
  \]

  \begin{itemize}
    \item Gradients are \emph{linear} in the parameters.
    \item Setting both to zero yields a system of linear equations: the \textbf{normal equations}.
    \item Exact, analytical solution — no iteration required.
  \end{itemize}

  \vspace{.4em}
  \begin{alertblock}{Limitation}
    A single large outlier can dominate the sum and distort the entire fitted model.
  \end{alertblock}
\end{frame}

% ---- 17 ----
\begin{frame}{Absolute Loss ($L_1$)}
  \begin{block}{Definition}
    \[
      L(y,t) = |y-t|
    \]
  \end{block}

  \begin{itemize}
    \item Errors penalised \emph{linearly}: doubling the error doubles the loss.
    \item Much more \textbf{robust to outliers} than quadratic loss.
    \item Gradient: piecewise constant, $\pm 1$ — steps of uniform size during GD.
    \item Convex but \emph{not} strictly convex; gradient undefined at $y=t$.
    \item Multiple optimal solutions possible; handled via \textbf{subgradient methods}.
  \end{itemize}
\end{frame}

% ---- 18 ----
\begin{frame}{Absolute Loss — Plot}
  \begin{center}
  \begin{tikzpicture}[scale=0.85]
    \begin{axis}[
      legend style={at={(1.5,1)},anchor=north east,draw=none,fill=none},
      ytick={-1,0,1}, yticklabels={$-1$,$0$,$1$},
      xticklabels={,,},
      axis x line=bottom, axis y line=left,
      xlabel={$y$},
      x label style={at={(axis description cs:0.99,0)},anchor=north},
      xmin=-2, xmax=2, ymin=-1.1, ymax=3.5,
      width=8cm, height=5cm]
      \addplot [cerise, thick, domain=-2:0, samples=201]{-x};
      \addplot [cerise, thick, domain=0:2, samples=201]{x};
      \addplot [slateblue, thick, domain=-2:0, samples=201]{-1};
      \addplot [slateblue, thick, domain=0:2, samples=201]{1};
      \addplot [dotted, domain=-3:3, samples=2] coordinates {(-3,0)(3,0)};
      \legend{loss, loss gradient}
    \end{axis}
  \end{tikzpicture}
  \end{center}
  {\small Absolute loss function and its gradient.
  The V-shape is characteristic, with constant gradient magnitude.}
\end{frame}

% ---- 19 ----
\begin{frame}{Huber Loss}
  \begin{block}{Definition}
    \[
      L(t,y)=
      \begin{cases}
        \tfrac{1}{2}(t-y)^2 & |t-y|\leq\delta \\[4pt]
        \delta\!\left(|t-y|-\tfrac{\delta}{2}\right) & |t-y|>\delta
      \end{cases}
    \]
  \end{block}

  \begin{itemize}
    \item \textbf{Quadratic} for small errors $\Rightarrow$ smooth gradients, efficient optimisation.
    \item \textbf{Linear} for large errors $\Rightarrow$ robust to outliers.
    \item $\delta$ is a hyperparameter controlling the transition point.
    \item Combines the best properties of $L_2$ and $L_1$.
  \end{itemize}
\end{frame}

% ---- 20 ----
\begin{frame}{Huber Loss — Plot}
  \begin{center}
  \begin{tikzpicture}[scale=0.85]
    \begin{axis}[
      legend style={at={(1.5,1)},anchor=north east,draw=none,fill=none},
      ytick={-1,0,1}, yticklabels={$-\delta$,$0$,$\delta$},
      xticklabels={,,},
      axis x line=bottom, axis y line=left,
      xlabel={$y$},
      x label style={at={(axis description cs:0.99,0)},anchor=north},
      xmin=-2, xmax=2, ymin=-1.1, ymax=3.5,
      width=8cm, height=5cm]
      \addplot [cerise, thick, domain=-1:1, samples=201]{x^2/2};
      \addplot [cerise, thick, domain=-2:-1, samples=201]{-x-1/2};
      \addplot [cerise, thick, domain=1:2, samples=201]{x-1/2};
      \addplot [slateblue, thick, domain=-1:1, samples=201]{x};
      \addplot [slateblue, thick, domain=-3:-1, samples=201]{-1};
      \addplot [slateblue, thick, domain=1:3, samples=201]{1};
      \addplot [dotted, domain=-3:3, samples=2] coordinates {(-3,0)(3,0)};
      \legend{loss, loss gradient}
    \end{axis}
  \end{tikzpicture}
  \end{center}
  {\small Huber loss combines quadratic behaviour near zero with linear growth for large errors.}
\end{frame}

% ---- 21 ----
\begin{frame}{Comparison: Regression Losses}
  \begin{center}
  \renewcommand{\arraystretch}{1.3}
  \begin{tabular}{lccc}
    \toprule
    & \textbf{$L_2$} & \textbf{$L_1$} & \textbf{Huber} \\
    \midrule
    Formula & $(y-t)^2$ & $|y-t|$ & piecewise \\
    Convexity & strictly & convex & strictly \\
    Smoothness & everywhere & not at $y=t$ & everywhere \\
    Outlier robustness & \alert{low} & high & high \\
    Gradient descent & direct & subgradient & direct \\
    Closed-form (linear) & \checkmark & $\times$ & $\times$ \\
    \bottomrule
  \end{tabular}
  \end{center}
\end{frame}

% ============================================================
\section{Loss Functions for Classification}
% ============================================================

% ---- 22 ----
\begin{frame}{Classification — Setup}
  \begin{itemize}
    \item Goal: assign inputs to \emph{discrete} classes.
    \item Focus on \textbf{binary classification}: $t \in \{-1, +1\}$.
    \item Classifier outputs a real value $y = h(\x)$:
      \begin{itemize}
        \item $\text{sign}(y)$ determines the predicted class.
        \item $|y|$ encodes prediction confidence.
      \end{itemize}
    \item Two approaches:
  \end{itemize}

  \vspace{.3em}
  \begin{columns}[T]
    \column{.49\textwidth}
    \begin{block}{Direct prediction}
      Returns a class label $\hat t$.\\
      Error if $\hat t \neq t$.
    \end{block}
    \column{.49\textwidth}
    \begin{block}{Probabilistic prediction}
      Returns $p(t|\x)$ for each class.\\
      Error measured via distribution divergence.
    \end{block}
  \end{columns}
\end{frame}

% ---- 23 ----
\begin{frame}{0/1 Loss}
  \begin{block}{Definition}
    \[
      L(t,y) = \mathbb{1}[ty < 0]
      \quad\Leftrightarrow\quad
      L(t,y)=\begin{cases} 1 & \text{sgn}(y)\neq t \\ 0 & \text{sgn}(y)=t \end{cases}
    \]
  \end{block}

  \begin{itemize}
    \item Directly minimises classification error rate. Intuitive.
    \item \alert{Critical problems}:
      \begin{itemize}
        \item \textbf{Non-convex}: many local minima; no optimality guarantee.
        \item \textbf{Non-smooth}: discontinuous jump at $ty=0$.
        \item \textbf{Zero gradient almost everywhere}: gradient descent cannot progress.
        \item \textbf{NP-hard}: minimising $\EmpR$ for linear classifiers is computationally intractable.
      \end{itemize}
  \end{itemize}
\end{frame}

% ---- 24 ----
\begin{frame}{0/1 Loss — Plot}
  \begin{center}
  \begin{tikzpicture}[scale=0.85]
    \begin{axis}[
      ytick={0,1}, yticklabels={$0$,$1$},
      xticklabels={,,},
      axis x line=bottom, axis y line=left,
      xlabel={$ty$},
      x label style={at={(axis description cs:0.99,0)},anchor=north},
      xmin=-2, xmax=2, ymin=-.1, ymax=3.5,
      width=8cm, height=5cm]
      \addplot [cerise, thick, domain=-2:0, samples=201]{1};
      \addplot [cerise, thick, domain=0:2, samples=201]{0};
      \addplot [dotted, domain=-3:3, samples=2] coordinates {(-3,0)(3,0)};
      \addplot [dotted, domain=-3:3, samples=2] coordinates {(-3,1)(3,1)};
    \end{axis}
  \end{tikzpicture}
  \end{center}
  {\small The 0/1 loss is a step function that assigns uniform cost to all errors.}
\end{frame}

% ---- 25 ----
\begin{frame}{Convex Surrogate Loss Functions}
  \begin{block}{Motivation}
    0/1 loss is ideal but intractable. We replace it with a \textbf{surrogate} that is:
  \end{block}

  \vspace{.3em}
  \begin{enumerate}
    \item \textbf{Upper bound} on 0/1 loss:
      $L_{\text{surr}}(t,y) \geq \mathbb{1}[ty<0]$
    \item \textbf{Convex} in $y$ (and in parameters).
    \item \textbf{Smooth}: well-defined derivatives almost everywhere.
  \end{enumerate}

  \vspace{.5em}
  \begin{exampleblock}{Key surrogates}
    Perceptron loss, Hinge loss, Log loss (cross-entropy), Squared loss (classification), Exponential loss.
  \end{exampleblock}

  \vspace{.3em}
  Minimising a surrogate $\Rightarrow$ tends to minimise classification error indirectly.
\end{frame}

% ---- 26 ----
\begin{frame}{Perceptron Loss}
  \begin{block}{Definition}
    \[
      L(t,y) = \max(0,\,-ty)
    \]
  \end{block}

  \begin{itemize}
    \item \textbf{Correctly classified} ($ty>0$): loss $= 0$. No penalty for correct predictions.
    \item \textbf{Misclassified} ($ty<0$): loss $= |ty|$. Severity grows with confidence of error.
    \item Continuous; gradient defined almost everywhere.
    \item Convex but \emph{not} strictly convex.
    \item \alert{Not a proper surrogate}: does not always upper-bound 0/1 loss.
    \item Historical significance: underlies the \textbf{classical Perceptron algorithm} (1950s).
  \end{itemize}
\end{frame}

% ---- 27 ----
\begin{frame}{Hinge Loss}
  \begin{block}{Definition}
    \[
      L(t,y) = \max(0,\,1-ty)
    \]
  \end{block}

  \begin{itemize}
    \item Loss $= 0$ only when $ty \geq 1$ (correct \emph{and} confident).
    \item $0 < ty < 1$: correct but not confident — still penalised.
    \item $ty < 0$: misclassified — loss grows linearly.
    \item Properties: continuous, convex, proper surrogate (upper-bounds 0/1 loss).
    \item \alert{Not differentiable} at $ty=1$ (corner); handled via \textbf{subgradients}.
    \item Foundation of \textbf{Support Vector Machines (SVM)}.
  \end{itemize}
\end{frame}

% ---- 28 ----
\begin{frame}{Hinge Loss — Plot}
  \begin{center}
  \begin{tikzpicture}[scale=0.85]
    \begin{axis}[
      legend style={at={(1.5,1)},anchor=north east,draw=none,fill=none},
      ytick={0,1}, yticklabels={$0$,$1$},
      xticklabels={,,},
      axis x line=bottom, axis y line=left,
      xlabel={$ty$},
      x label style={at={(axis description cs:0.99,0)},anchor=north},
      xmin=-2, xmax=2, ymin=-.1, ymax=3.5,
      width=8cm, height=5cm]
      \addplot [cerise, thick, domain=1:2, samples=201]{0};
      \addplot [cerise, thick, domain=-2:1, samples=201]{1-x};
      \addplot [teal, thick, domain=-2:0, samples=201]{1};
      \addplot [teal, thick, domain=0:2, samples=201]{0};
      \addplot [dotted, domain=-3:3, samples=2] coordinates {(-3,0)(3,0)};
      \addplot [dotted, domain=-3:3, samples=2] coordinates {(-3,1)(3,1)};
      \legend{Hinge loss, 0/1 loss}
    \end{axis}
  \end{tikzpicture}
  \end{center}
\end{frame}

% ============================================================
\section{Subgradients}
% ============================================================

% ---- 29 ----
\begin{frame}{Non-Differentiable Losses — The Problem}
  Hinge loss: $L_H(t,y)=\max(0,1-ty)$

  \[
    \frac{\partial}{\partial y}L_H =
    \begin{cases}
      -t & ty < 1 \\
      0  & ty > 1 \\
      \text{undefined} & ty = 1
    \end{cases}
  \]

  \begin{itemize}
    \item The gradient is \alert{undefined} at the corner $ty=1$.
    \item Standard gradient descent cannot be applied directly.
    \item Same issue arises for absolute loss at $y=t$ and perceptron loss at $ty=0$.
  \end{itemize}

  \vspace{.4em}
  \begin{block}{Solution: subgradients}
    Generalise the gradient to non-smooth points using a \emph{set} of valid linear bounds.
  \end{block}
\end{frame}

% ---- 30 ----
\begin{frame}{Subgradients — Definition}
  For a convex function $f$, at a differentiable point $x$:
  \[
    f(x') \geq f(x) + \nabla f|_x \cdot (x'-x) \quad \forall\, x'
  \]
  (Gradient defines a tangent lower bound.)

  \vspace{.5em}
  \begin{block}{Subgradient}
    A vector $\bar\nabla f|_x$ is a \textbf{subgradient} at $x$ if:
    \[
      f(x') \geq f(x) + \bar\nabla f|_x \cdot (x'-x) \quad \forall\, x'
    \]
  \end{block}

  \begin{itemize}
    \item At differentiable points: the subgradient is unique and equals the gradient.
    \item At non-differentiable points: there is a \emph{set} of valid subgradients.
    \item Any element of this set can be used in place of the gradient in GD.
  \end{itemize}
\end{frame}

% ---- 31 ----
\begin{frame}{Subgradient of Hinge Loss}
  At the corner $ty=1$, valid slopes lie between $-t$ and $0$.

  \vspace{.3em}
  \begin{center}
  \begin{tikzpicture}[scale=0.85]
    \begin{axis}[
      legend style={at={(1.5,1)},anchor=north east,draw=none,fill=none},
      ytick={0,1}, yticklabels={$0$,$1$},
      xticklabels={,,},
      axis x line=bottom, axis y line=left,
      xlabel={$ty$},
      x label style={at={(axis description cs:0.99,0)},anchor=north},
      xmin=-2, xmax=3, ymin=-.9, ymax=3.5,
      width=8cm, height=4.5cm]
      \addplot [cerise, thick, domain=1:3]{0};
      \addplot [cerise, thick, domain=-2:1]{1-x};
      \addplot [orange, dashed, domain=-2:3]{.5-.5*x};
      \addplot [slateblue, dashed, domain=-2:3]{.3-.3*x};
      \addplot [dotted, domain=-3:4] coordinates {(-3,0)(4,0)};
      \addplot [dotted, domain=-3:4] coordinates {(-3,1)(4,1)};
    \end{axis}
  \end{tikzpicture}
  \end{center}

  A common choice: set the subgradient to $0$ at $ty=1$, giving:
  \[
    \bar\nabla_y L_H =
    \begin{cases}
      -t & ty < 1 \\
      0  & ty \geq 1
    \end{cases}
  \]
\end{frame}

% ---- 32 ----
\begin{frame}{Squared Loss for Classification}
  \begin{block}{Definition}
    \[
      L(t,y) = (1-ty)^2
    \]
  \end{block}

  \begin{columns}[T]
    \column{.5\textwidth}
    \begin{block}{Advantages}
      \begin{itemize}
        \item Infinitely differentiable.
        \item Convex, well-defined gradient.
        \item No corners or undefined points.
      \end{itemize}
    \end{block}
    \column{.5\textwidth}
    \begin{alertblock}{Disadvantages}
      \begin{itemize}
        \item \emph{Not} a proper surrogate: does not upper-bound 0/1 loss.
        \item Sensitive to outliers (quadratic growth).
        \item Penalises \emph{very confident correct} predictions.
      \end{itemize}
    \end{alertblock}
  \end{columns}

  \vspace{.4em}
  \begin{itemize}
    \item Rarely the preferred choice in classification practice.
  \end{itemize}
\end{frame}

% ---- 33 ----
\begin{frame}{Squared Loss for Classification — Plot}
  \begin{center}
  \begin{tikzpicture}[scale=0.85]
    \begin{axis}[
      legend style={at={(1.5,1)},anchor=north east,draw=none,fill=none},
      ytick={0,1}, yticklabels={$0$,$1$},
      xticklabels={,,},
      axis x line=bottom, axis y line=left,
      xlabel={$ty$},
      x label style={at={(axis description cs:0.99,0)},anchor=north},
      xmin=-2, xmax=2, ymin=-.1, ymax=3.5,
      width=8cm, height=5cm]
      \addplot [cerise, thick, domain=-2:2]{(1-x)^2};
      \addplot [teal, thick, domain=-2:0]{1};
      \addplot [teal, thick, domain=0:2]{0};
      \addplot [dotted, domain=-3:3] coordinates {(-3,0)(3,0)};
      \addplot [dotted, domain=-3:3] coordinates {(-3,1)(3,1)};
      \legend{Square loss, 0/1 loss}
    \end{axis}
  \end{tikzpicture}
  \end{center}
\end{frame}

% ============================================================
\section{Log Loss and Information Theory}
% ============================================================

% ---- 34 ----
\begin{frame}{Log Loss (Cross-Entropy Loss)}
  \begin{block}{Definition}
    \[
      L(t,y) = \frac{1}{\log 2}\log\!\left(1+e^{-ty}\right)
    \]
  \end{block}

  \begin{itemize}
    \item Also called \textbf{logistic loss} or \textbf{cross-entropy loss}.
    \item Standard loss for probabilistic classification and logistic regression.
    \item Smooth approximation to hinge loss (no corner).
    \item For large negative $ty$: grows approximately linearly.
    \item For large positive $ty$: approaches 0 asymptotically (never reaches it).
  \end{itemize}
\end{frame}

% ---- 35 ----
\begin{frame}{Log Loss — Properties}
  \begin{columns}[T]
    \column{.55\textwidth}
    \begin{block}{Mathematical properties}
      \begin{itemize}
        \item Infinitely differentiable everywhere.
        \item \textbf{Strictly convex}: unique global minimum.
        \item Proper surrogate: upper-bounds 0/1 loss.
        \item Gradient continuous and well-defined.
      \end{itemize}
    \end{block}
    \column{.45\textwidth}
    \begin{alertblock}{Caution}
      \begin{itemize}
        \item Grows unboundedly for $ty\to-\infty$.
        \item Sensitive to severely mislabelled examples.
        \item Growth is \emph{logarithmic}, so more robust than squared loss.
      \end{itemize}
    \end{alertblock}
  \end{columns}
\end{frame}

% ---- 36 ----
\begin{frame}{Log Loss — Plot}
  \begin{center}
  \begin{tikzpicture}[scale=0.85]
    \begin{axis}[
      legend style={at={(1.5,1)},anchor=north east,draw=none,fill=none},
      ytick={0,1}, yticklabels={$0$,$1$},
      xticklabels={,,},
      axis x line=bottom, axis y line=left,
      xlabel={$ty$},
      x label style={at={(axis description cs:0.99,0)},anchor=north},
      xmin=-2, xmax=2, ymin=-.1, ymax=3.5,
      width=8cm, height=5cm]
      \addplot [cerise, thick, domain=-2:2, samples=200]{ln(1+exp(-x))/ln(2)};
      \addplot [teal, thick, domain=-2:0]{1};
      \addplot [teal, thick, domain=0:2]{0};
      \addplot [dotted, domain=-3:3] coordinates {(-3,0)(3,0)};
      \addplot [dotted, domain=-3:3] coordinates {(-3,1)(3,1)};
      \legend{Log loss, 0/1 loss}
    \end{axis}
  \end{tikzpicture}
  \end{center}
\end{frame}

% ---- 37 ----
\begin{frame}{Information Theory Background}
  \begin{columns}[T]
    \column{.33\textwidth}
    \begin{block}{Entropy $H(p)$}
      \[
        H(p) = -\mathbb{E}_{x\sim p}[\log p(x)]
      \]
      Uncertainty / information content of $p$.\\
      Min.\ avg.\ bits to encode from $p$.
    \end{block}
    \column{.33\textwidth}
    \begin{block}{Cross-entropy $\mathrm{H}(p,q)$}
      \[
        \mathrm{H}(p,q) = -\mathbb{E}_{x\sim p}[\log q(x)]
      \]
      Avg.\ bits using code for $q$, when true dist.\ is $p$.
    \end{block}
    \column{.33\textwidth}
    \begin{block}{KL divergence}
      \[
        \mathrm{KL}(p\|q) = \mathrm{H}(p,q)-H(p)
      \]
      Extra bits wasted by using $q$ instead of $p$.\\
      Always $\geq 0$; $=0$ iff $p=q$.
    \end{block}
  \end{columns}

  \vspace{.5em}
  \begin{itemize}
    \item $\mathrm{KL}(p\|q)=0 \Leftrightarrow p=q$ (almost everywhere).
    \item Minimising cross-entropy $\equiv$ minimising KL divergence (since $H(p)$ is fixed).
  \end{itemize}
\end{frame}

% ---- 38 ----
\begin{frame}{Cross-Entropy for Binary Classification}
  Let $p(\x)\in\{0,1\}$ be the true label, $y(\x)\in[0,1]$ the predicted probability for class $C_1$.

  \begin{block}{Empirical cross-entropy}
    \[
      \mathrm{CE}(\calT)
      = -\frac{1}{|\calT|}\sum_{(\x,t)\in\calT}
        \Big[t\log y(\x) + (1-t)\log(1-y(\x))\Big]
    \]
  \end{block}

  \begin{itemize}
    \item For each example: $\log$ probability assigned to the \emph{correct} class, averaged.
    \item The negative sign converts high probability (good) into low loss.
  \end{itemize}

  \vspace{.3em}
  Separating by class:
  \[
    \mathrm{CE}(\calT)
    = -\frac{1}{|\calT|}
      \left(\sum_{(\x,t)\in C_1}\!\!\log y(\x)
            +\sum_{(\x,t)\in C_0}\!\!\log(1-y(\x))\right)
  \]
\end{frame}

% ---- 39 ----
\begin{frame}{Log Loss $\equiv$ Cross-Entropy (Logistic Regression)}
  For logistic regression: $y(\x)=\sigma(\w^T\x+b)=\dfrac{1}{1+e^{-(\w^T\x+b)}}$

  \vspace{.5em}
  Substituting into cross-entropy, using $t\in\{-1,+1\}$:
  \[
    \mathrm{CE}(\calT)
    = \frac{1}{|\calT|}\sum_{(\x,t)\in\calT}\log\!\left(1+e^{-t(\w^T\x+b)}\right)
  \]

  \begin{block}{Connection to log loss}
    \[
      \EmpR_{\calT}(h)
      = \frac{1}{|\calT|\log 2}\sum_{(\x,t)\in\calT}\log\!\left(1+e^{-t(\w^T\x+b)}\right)
    \]
  \end{block}

  \begin{itemize}
    \item Minimising log loss $\equiv$ minimising cross-entropy.
    \item Strong information-theoretic justification for using log loss in classification.
  \end{itemize}
\end{frame}

% ============================================================
\section{Exponential Loss}
% ============================================================

% ---- 40 ----
\begin{frame}{Exponential Loss}
  \begin{block}{Definition}
    \[
      L(t,y) = e^{-ty}
    \]
  \end{block}

  \begin{itemize}
    \item Prominent in \textbf{boosting} algorithms, especially \textbf{AdaBoost}.
    \item Penalises wrong predictions more aggressively than log loss: exponential growth.
    \item Infinitely differentiable everywhere.
    \item Strictly convex $\Rightarrow$ unique global minimum.
    \item Proper surrogate: upper-bounds 0/1 loss.
    \item \alert{Sensitive to outliers}: very large penalty for confident misclassifications.
  \end{itemize}
\end{frame}

% ---- 41 ----
\begin{frame}{Exponential Loss — Plot}
  \begin{center}
  \begin{tikzpicture}[scale=0.85]
    \begin{axis}[
      legend style={at={(1.5,1)},anchor=north east,draw=none,fill=none},
      ytick={0,1}, yticklabels={$0$,$1$},
      xticklabels={,,},
      axis x line=bottom, axis y line=left,
      xlabel={$ty$},
      x label style={at={(axis description cs:0.99,0)},anchor=north},
      xmin=-2, xmax=2, ymin=-.1, ymax=3.5,
      width=8cm, height=5cm]
      \addplot [cerise, thick, domain=-2:2, samples=200]{exp(-x)};
      \addplot [teal, thick, domain=-2:0]{1};
      \addplot [teal, thick, domain=0:2]{0};
      \addplot [dotted, domain=-3:3] coordinates {(-3,0)(3,0)};
      \addplot [dotted, domain=-3:3] coordinates {(-3,1)(3,1)};
      \legend{Exponential loss, 0/1 loss}
    \end{axis}
  \end{tikzpicture}
  \end{center}
\end{frame}

% ---- 42 ----
\begin{frame}{Comparison: Classification Losses}
  \scriptsize
  \renewcommand{\arraystretch}{1.35}
  \begin{center}
  \begin{tabular}{l|ccccc}
    \toprule
    & \textbf{0/1} & \textbf{Perceptron} & \textbf{Hinge} & \textbf{Log} & \textbf{Exp.} \\
    \midrule
    Upper-bounds 0/1 & --- & partial & \checkmark & \checkmark & \checkmark \\
    Convex           & \alert{No}  & Yes & Yes & Yes & Yes \\
    Strictly convex  & No  & No  & No  & \checkmark & \checkmark \\
    Smooth           & No  & No  & No  & \checkmark & \checkmark \\
    Outlier robust   & ---  & medium & medium & medium & \alert{low} \\
    Use case         & ideal & Perceptron & SVM & Logistic Reg. & AdaBoost \\
    \bottomrule
  \end{tabular}
  \end{center}
\end{frame}

% ============================================================
\section{Stochastic Gradient Descent}
% ============================================================

% ---- 43 ----
\begin{frame}{Stochastic Gradient Descent (SGD)}
  \begin{block}{Motivation}
    Batch GD computes the gradient over the \emph{entire} training set at each step.
    This is expensive when $|\calT|$ is large.
  \end{block}

  \vspace{.4em}
  \begin{block}{SGD idea}
    Use a \alert{single randomly chosen example} $(\x_i, t_i)$ to estimate the gradient:
    \[
      \Vtheta^{(i)} \leftarrow \Vtheta^{(i-1)}
      - \eta\,\nabla L\!\left(h_{\Vtheta}(\x_i), t_i\right)\Big|_{\Vtheta^{(i-1)}}
    \]
  \end{block}

  \begin{itemize}
    \item Noisy gradient estimate, but very fast per step.
    \item Noise can \emph{help} escape local minima.
    \item Widely used in deep learning practice.
  \end{itemize}
\end{frame}

% ---- 44 ----
\begin{frame}{Mini-Batch Gradient Descent}
  \begin{block}{Compromise}
    Use a \textbf{mini-batch} $\mathcal{B} \subset \calT$ of $m$ randomly chosen examples:
    \[
      \Vtheta^{(i)} \leftarrow \Vtheta^{(i-1)}
      - \frac{\eta}{m}\sum_{(\x,t)\in\mathcal{B}}
        \nabla L\!\left(h_{\Vtheta}(\x), t\right)\Big|_{\Vtheta^{(i-1)}}
    \]
  \end{block}

  \vspace{.4em}
  \begin{center}
  \renewcommand{\arraystretch}{1.2}
  \begin{tabular}{lccc}
    \toprule
    & \textbf{Batch} & \textbf{Mini-batch} & \textbf{SGD} \\
    \midrule
    Gradient quality & exact & good estimate & noisy \\
    Speed per step   & slow  & moderate & fast \\
    Memory use       & full  & moderate & minimal \\
    Typical use      & small $n$ & deep learning & large $n$ \\
    \bottomrule
  \end{tabular}
  \end{center}
\end{frame}

% ---- 45 ----
\begin{frame}{SGD — Key Properties}
  \begin{columns}[T]
    \column{.5\textwidth}
    \begin{block}{Advantages}
      \begin{itemize}
        \item Each update is $O(d)$ — independent of $n$.
        \item Can process data in a streaming fashion.
        \item Noise helps escape poor local minima.
        \item Enables learning from very large datasets.
      \end{itemize}
    \end{block}
    \column{.5\textwidth}
    \begin{alertblock}{Challenges}
      \begin{itemize}
        \item High variance in gradient estimates.
        \item May not converge to exact minimum.
        \item Learning rate tuning more critical.
        \item Convergence guarantees are weaker.
      \end{itemize}
    \end{alertblock}
  \end{columns}

  \vspace{.5em}
  \begin{exampleblock}{Modern practice}
    Almost all deep learning uses mini-batch SGD with adaptive learning rates
    (Adam, RMSProp, AdaGrad).
  \end{exampleblock}
\end{frame}

% ============================================================
\section{Summary}
% ============================================================

% ---- 46 ----
\begin{frame}{Summary — Loss Functions}
  \begin{block}{Core definitions}
    \begin{itemize}
      \item \textbf{Loss} $L(y,t)$: pointwise prediction penalty.
      \item \textbf{Empirical risk} $\EmpR_\calT(h)$: average loss over training set.
      \item \textbf{Training} $\equiv$ minimising $\EmpR_\calT(h_{\Vtheta})$ over parameters $\Vtheta$.
    \end{itemize}
  \end{block}

  \vspace{.4em}
  \begin{block}{Regression losses}
    \begin{itemize}
      \item Quadratic ($L_2$): smooth, convex, sensitive to outliers.
      \item Absolute ($L_1$): robust, requires subgradients.
      \item Huber: best of both; controlled by $\delta$.
    \end{itemize}
  \end{block}

  \vspace{.4em}
  \begin{block}{Classification losses}
    \begin{itemize}
      \item 0/1: intuitive but NP-hard to optimise.
      \item Surrogates (Hinge, Log, Exponential): convex, tractable.
      \item Log loss $\equiv$ cross-entropy: information-theoretic grounding.
    \end{itemize}
  \end{block}
\end{frame}

% ---- 47 ----
%\begin{frame}{Summary — Optimisation}
%  \begin{block}{Gradient descent}
%    \begin{itemize}
%      \item Iteratively follow the negative gradient.
%      \item Learning rate $\eta$ controls step size.
%      \item Converges to a local minimum (global if loss is convex).
%    \end{itemize}
%  \end{block}
%
%  \vspace{.4em}
%  \begin{block}{Convexity}
%    \begin{itemize}
%      \item Convex loss: every local minimum is global.
%      \item Strictly convex: unique global minimum.
%      \item Sum of convex functions is convex $\Rightarrow$ empirical risk inherits convexity.
%    \end{itemize}
%  \end{block}
%
%  \vspace{.4em}
%  \begin{block}{Variants}
%    \begin{itemize}
%      \item Mini-batch / SGD: practical for large-scale learning.
%      \item Subgradients: extend GD to non-smooth losses.
%    \end{itemize}
%  \end{block}
%\end{frame}

% ---- 48 — final message ----
\begin{frame}{Central Message}
  \begin{center}
    \Large
    \textbf{The loss function encodes what ``good predictions'' means.}
  \end{center}

  \vspace{.5em}
  \begin{itemize}
    \item Choosing the right loss is as important as choosing the model.
    \item \textbf{Regression}: prefer $L_2$ for smooth data, $L_1$/Huber for outliers.
    \item \textbf{Classification}: use convex surrogates; log loss has the strongest theoretical backing.
    \item \textbf{Optimisation}: exploit convexity when possible; use SGD variants in practice.
    \item The interplay of loss, optimiser, and model defines the entire learning procedure.
  \end{itemize}

%  \vspace{.8em}
%  \begin{center}
%    \textcolor{torblue}{\rule{.6\textwidth}{.4pt}}\\[.3em]
%    {\small Giorgio Gambosi — Foundations of Machine Learning — A.Y.\ 2025–2026}
%  \end{center}
\end{frame}

%\end{document}
